% page 124 du fac-simile (image p-123 du scan)
% bloc 160.0 x 233.3 mm ; marge gauche 8.0 mm
\pgc{-12.700mm}{0.000mm}{
\l{\cel{7}114.}
\l{}
\l{\sou{diskursar}. (\sou{netrans., pri)}. I. Developar temo oratorala ko-}
\l{\cel{3}ram un o plura personi, animante la pronunco per gesti.}
\l{\cel{3}- II. Developar docive temo (parole o skribe). - DEFIRS.}
\l{}
\l{\sou{diskursiva}. (\sou{logiko}) Qua pasas de ideo ad altra, qua pro-}
\l{\cel{3}cedas per serchi inter-sucedanta (opoze ad intuicala). -}
\l{\cel{3}DEF.}
\l{}
\l{\sou{diskutar}. (\sou{trans}.) Studiar problemo, questiono, ponderante}
\l{\cel{3}la \textquotedbl{}por\textquotedbl{} e la \textquotedbl{}kontre\textquotedbl{}. (Aludante\cel{1}\sou{personi qui ne havas}}
\l{\cel{3}la\cel{1}\sou{sama opiniono)}. Kambiar argumento-punti pri ula temo.}
\l{\cel{3}- DEFIRS.}
\l{}
\l{\sou{dislokar}. (\sou{trans.}) Separar violentoze (la parti di ensemblo).}
\l{\cel{3}- DEFIS.}
\l{}
\l{\sou{dismo}. I. Dekimo de la rekoltajo (che la Judi, olim) presus-}
\l{\cel{3}tracionata por ofrar a la Sinioro o donata a la levidi.}
\l{\cel{3}- II. Parto (variiva) de la rekoltajo, presustracionata da}
\l{\cel{3}la eklezio katolika e da la siniori feudala. - FIS.}
\l{}
\l{\sou{disociar}.- (\sou{trans.}) (\sou{kemio)}. Separar la elementi unionita.}
\l{\cel{3}(\sou{anke metaf.}). - DEFIRS.}
\l{}
\l{\sou{disonancar}. (\sou{netrans.})\cel{2}Formacar harmonio poke agreabla}
\l{\cel{3}audale. - DEFIRS.}
\l{}
\l{\sou{dispacho}. (\sou{komerco}). Aranjo, konvenciono pri la perdi, la ava-}
\l{\cel{3}rii, inter kompanio di asekuro navala e la asekurito. - DFS.}
\l{}
\l{\sou{disparata}. Qua shokante dessimilesas to quo cirkondas lu. -}
\l{\cel{3}DEFIS.}
\l{}
\l{\sou{dispensar}.\cel{2}(\sou{trans.,}\cel{1}de).\cel{2}I. Permisar specale (ad ulu) agar}
\l{\cel{3}o facar ulo quon interdiktas la regulo di la eklezio kato-}
\l{\cel{3}lika, o ne agar o facar to quon lu preskriptas. - II. Per-}
\l{\cel{3}misar (ad ulu) ne agar o facar ulo. - DEFIRS.}
\l{}
\l{\sou{disp}ensario. Loko en qua la povri povas konsultar mediki e}
\l{\cel{3}recevar flegi o medikamenti (sen lojar en olu). - EFIS.}
\l{}
\l{\sou{dispepsio}.\cel{2}Desfacileso pri digestar. - DEFIS.}
\l{}
\l{\sou{dispersar}. (\sou{trans}.) Separar per pulsar li, per igar li dis-}
\l{\cel{3}irar (kozi, personi, qui asemblesis). - DEFIS.}
\l{}
\l{\sou{dispiremo}. (\sou{histologio}) Spiremo duopla qua, dum mitoso, apa-}
\l{\cel{3}ras an la poli di la celulo, kande la kromosomi rigrupigas}
\l{\cel{3}su por formacar la elementi nuklea di la du celuli nova.}
\l{}
\l{disponar. (trans.) Agar segun sua arbitrio pri ulo, pri ulu;}
\l{\cel{3}darfar agar, uzar, ulo quale onu volas lo. - DEFIRS.}
}
